% Options for packages loaded elsewhere
\PassOptionsToPackage{unicode}{hyperref}
\PassOptionsToPackage{hyphens}{url}
%
\documentclass[
  ignorenonframetext,
]{beamer}
\usepackage{pgfpages}
\setbeamertemplate{caption}[numbered]
\setbeamertemplate{caption label separator}{: }
\setbeamercolor{caption name}{fg=normal text.fg}
\beamertemplatenavigationsymbolsempty
% Prevent slide breaks in the middle of a paragraph
\widowpenalties 1 10000
\raggedbottom
\setbeamertemplate{part page}{
  \centering
  \begin{beamercolorbox}[sep=16pt,center]{part title}
    \usebeamerfont{part title}\insertpart\par
  \end{beamercolorbox}
}
\setbeamertemplate{section page}{
  \centering
  \begin{beamercolorbox}[sep=12pt,center]{part title}
    \usebeamerfont{section title}\insertsection\par
  \end{beamercolorbox}
}
\setbeamertemplate{subsection page}{
  \centering
  \begin{beamercolorbox}[sep=8pt,center]{part title}
    \usebeamerfont{subsection title}\insertsubsection\par
  \end{beamercolorbox}
}
\AtBeginPart{
  \frame{\partpage}
}
\AtBeginSection{
  \ifbibliography
  \else
    \frame{\sectionpage}
  \fi
}
\AtBeginSubsection{
  \frame{\subsectionpage}
}
\usepackage{amsmath,amssymb}
\usepackage{lmodern}
\usepackage{iftex}
\ifPDFTeX
  \usepackage[T1]{fontenc}
  \usepackage[utf8]{inputenc}
  \usepackage{textcomp} % provide euro and other symbols
\else % if luatex or xetex
  \usepackage{unicode-math}
  \defaultfontfeatures{Scale=MatchLowercase}
  \defaultfontfeatures[\rmfamily]{Ligatures=TeX,Scale=1}
\fi
\usetheme[]{Montpellier}
% Use upquote if available, for straight quotes in verbatim environments
\IfFileExists{upquote.sty}{\usepackage{upquote}}{}
\IfFileExists{microtype.sty}{% use microtype if available
  \usepackage[]{microtype}
  \UseMicrotypeSet[protrusion]{basicmath} % disable protrusion for tt fonts
}{}
\makeatletter
\@ifundefined{KOMAClassName}{% if non-KOMA class
  \IfFileExists{parskip.sty}{%
    \usepackage{parskip}
  }{% else
    \setlength{\parindent}{0pt}
    \setlength{\parskip}{6pt plus 2pt minus 1pt}}
}{% if KOMA class
  \KOMAoptions{parskip=half}}
\makeatother
\usepackage{xcolor}
\IfFileExists{xurl.sty}{\usepackage{xurl}}{} % add URL line breaks if available
\IfFileExists{bookmark.sty}{\usepackage{bookmark}}{\usepackage{hyperref}}
\hypersetup{
  hidelinks,
  pdfcreator={LaTeX via pandoc}}
\urlstyle{same} % disable monospaced font for URLs
\newif\ifbibliography
\usepackage{color}
\usepackage{fancyvrb}
\newcommand{\VerbBar}{|}
\newcommand{\VERB}{\Verb[commandchars=\\\{\}]}
\DefineVerbatimEnvironment{Highlighting}{Verbatim}{commandchars=\\\{\}}
% Add ',fontsize=\small' for more characters per line
\usepackage{framed}
\definecolor{shadecolor}{RGB}{248,248,248}
\newenvironment{Shaded}{\begin{snugshade}}{\end{snugshade}}
\newcommand{\AlertTok}[1]{\textcolor[rgb]{0.94,0.16,0.16}{#1}}
\newcommand{\AnnotationTok}[1]{\textcolor[rgb]{0.56,0.35,0.01}{\textbf{\textit{#1}}}}
\newcommand{\AttributeTok}[1]{\textcolor[rgb]{0.77,0.63,0.00}{#1}}
\newcommand{\BaseNTok}[1]{\textcolor[rgb]{0.00,0.00,0.81}{#1}}
\newcommand{\BuiltInTok}[1]{#1}
\newcommand{\CharTok}[1]{\textcolor[rgb]{0.31,0.60,0.02}{#1}}
\newcommand{\CommentTok}[1]{\textcolor[rgb]{0.56,0.35,0.01}{\textit{#1}}}
\newcommand{\CommentVarTok}[1]{\textcolor[rgb]{0.56,0.35,0.01}{\textbf{\textit{#1}}}}
\newcommand{\ConstantTok}[1]{\textcolor[rgb]{0.00,0.00,0.00}{#1}}
\newcommand{\ControlFlowTok}[1]{\textcolor[rgb]{0.13,0.29,0.53}{\textbf{#1}}}
\newcommand{\DataTypeTok}[1]{\textcolor[rgb]{0.13,0.29,0.53}{#1}}
\newcommand{\DecValTok}[1]{\textcolor[rgb]{0.00,0.00,0.81}{#1}}
\newcommand{\DocumentationTok}[1]{\textcolor[rgb]{0.56,0.35,0.01}{\textbf{\textit{#1}}}}
\newcommand{\ErrorTok}[1]{\textcolor[rgb]{0.64,0.00,0.00}{\textbf{#1}}}
\newcommand{\ExtensionTok}[1]{#1}
\newcommand{\FloatTok}[1]{\textcolor[rgb]{0.00,0.00,0.81}{#1}}
\newcommand{\FunctionTok}[1]{\textcolor[rgb]{0.00,0.00,0.00}{#1}}
\newcommand{\ImportTok}[1]{#1}
\newcommand{\InformationTok}[1]{\textcolor[rgb]{0.56,0.35,0.01}{\textbf{\textit{#1}}}}
\newcommand{\KeywordTok}[1]{\textcolor[rgb]{0.13,0.29,0.53}{\textbf{#1}}}
\newcommand{\NormalTok}[1]{#1}
\newcommand{\OperatorTok}[1]{\textcolor[rgb]{0.81,0.36,0.00}{\textbf{#1}}}
\newcommand{\OtherTok}[1]{\textcolor[rgb]{0.56,0.35,0.01}{#1}}
\newcommand{\PreprocessorTok}[1]{\textcolor[rgb]{0.56,0.35,0.01}{\textit{#1}}}
\newcommand{\RegionMarkerTok}[1]{#1}
\newcommand{\SpecialCharTok}[1]{\textcolor[rgb]{0.00,0.00,0.00}{#1}}
\newcommand{\SpecialStringTok}[1]{\textcolor[rgb]{0.31,0.60,0.02}{#1}}
\newcommand{\StringTok}[1]{\textcolor[rgb]{0.31,0.60,0.02}{#1}}
\newcommand{\VariableTok}[1]{\textcolor[rgb]{0.00,0.00,0.00}{#1}}
\newcommand{\VerbatimStringTok}[1]{\textcolor[rgb]{0.31,0.60,0.02}{#1}}
\newcommand{\WarningTok}[1]{\textcolor[rgb]{0.56,0.35,0.01}{\textbf{\textit{#1}}}}
\usepackage{longtable,booktabs,array}
\usepackage{calc} % for calculating minipage widths
\usepackage{caption}
% Make caption package work with longtable
\makeatletter
\def\fnum@table{\tablename~\thetable}
\makeatother
\setlength{\emergencystretch}{3em} % prevent overfull lines
\providecommand{\tightlist}{%
  \setlength{\itemsep}{0pt}\setlength{\parskip}{0pt}}
\setcounter{secnumdepth}{-\maxdimen} % remove section numbering
%\documentclass[unicode,10pt]{beamer}% 'unicode'が必要
\usepackage{luatexja}% 日本語を使う場合は必須
%\usepackage[japanese]{babel}	
%\usefonttheme[onlymath]{serif}
%\usepackage[T1]{fontenc} %これを使うと、ギリシャ文字が出ない
%\usepackage{textcomp}
%\usepackage[scale=1.0]{tgheros} % Sans serif font
%\usepackage[scaled]{beramono} % Monospace font
%\usepackage{luatexja-otf}
%\renewcommand{\kanjifamilydefault}{\gtdefault}
%\usepackage[haranoaji]{luatexja-preset}

% スライド番号の表示 %%%%%%%%%%%%%%%%%%%
\makeatletter
\setbeamertemplate{footline}{%
  \leavevmode%
  \hbox{%
  \begin{beamercolorbox}[wd=.5\paperwidth,ht=2.25ex,dp=1ex,right]{date in head/foot}%
    \insertframenumber{} \hspace*{2ex} % new version without total frames
  \end{beamercolorbox}}%
  \vskip0pt%
}
\makeatother
% スライド番号の表示 %%%%%%%%%%%%%%%%%%%
\ifLuaTeX
  \usepackage{selnolig}  % disable illegal ligatures
\fi

\title{第6章: 重回帰分析：発展}
\subtitle{Jeffrey Wooldridge (2016).\\
Introductory Econometrics: A Modern Approach\\
Seventh Edition. Cengage Learning.}
\author{}
\date{\vspace{-2.5em}2026-03-06}

\begin{document}
\frame{\titlepage}

\hypertarget{ux6e96ux5099}{%
\section{準備}\label{ux6e96ux5099}}

\begin{frame}[fragile]{必要なパッケージの読み込み}
\protect\hypertarget{ux5fc5ux8981ux306aux30d1ux30c3ux30b1ux30fcux30b8ux306eux8aadux307fux8fbcux307f}{}
\begin{itemize}
\tightlist
\item
  \texttt{wooldridge}パッケージの読み込み
\end{itemize}

\begin{Shaded}
\begin{Highlighting}[]
\FunctionTok{library}\NormalTok{(wooldridge)}
\end{Highlighting}
\end{Shaded}
\end{frame}

\hypertarget{ux30c7ux30fcux30bfux306eux30b9ux30b1ux30fcux30eaux30f3ux30b0ux3068olsux7d71ux8a08ux91cf}{%
\section{6-1
データのスケーリングとOLS統計量}\label{ux30c7ux30fcux30bfux306eux30b9ux30b1ux30fcux30eaux30f3ux30b0ux3068olsux7d71ux8a08ux91cf}}

\begin{frame}{単位の変更の影響}
\protect\hypertarget{ux5358ux4f4dux306eux5909ux66f4ux306eux5f71ux97ff}{}
\begin{itemize}
\tightlist
\item
  従属変数 \(y\) や独立変数 \(x\)
  の単位を変更（スケーリング）した場合、OLS推定量や統計量はどのように変化するか？
\item
  \textbf{例：} \(y\) を円から千円にする（1000で割る）、\(x\)
  をマイルからキロメートルにするなど。
\end{itemize}
\end{frame}

\begin{frame}{従属変数 \(y\) のスケーリング}
\protect\hypertarget{ux5f93ux5c5eux5909ux6570-y-ux306eux30b9ux30b1ux30fcux30eaux30f3ux30b0}{}
\begin{itemize}
\tightlist
\item
  \(y\) を定数 \(c\) 倍（例えば \(y^* = cy\)）した場合：

  \begin{itemize}
  \tightlist
  \item
    係数 \(\hat{\beta}_j\) はすべて \(c\) 倍になる。
  \item
    標準誤差 \(\text{se}(\hat{\beta}_j)\) も \(c\) 倍になる。
  \item
    \textbf{t統計量は変化しない。}
  \item
    \textbf{決定係数 \(R^2\) は変化しない。}
  \end{itemize}
\end{itemize}
\end{frame}

\begin{frame}{独立変数 \(x_j\) のスケーリング}
\protect\hypertarget{ux72ecux7acbux5909ux6570-x_j-ux306eux30b9ux30b1ux30fcux30eaux30f3ux30b0}{}
\begin{itemize}
\tightlist
\item
  独立変数 \(x_k\) を定数 \(c\) 倍（例えば \(x_k^* = cx_k\)）した場合：

  \begin{itemize}
  \tightlist
  \item
    該当する係数 \(\hat{\beta}_k\) は \(1/c\) 倍になる。
  \item
    他の係数 \(\hat{\beta}_j\) (\(j \neq k\)) は変化しない。
  \item
    標準誤差も同様に変化するため、\textbf{t統計量は変化しない。}
  \item
    \textbf{決定係数 \(R^2\) は変化しない。}
  \end{itemize}
\end{itemize}
\end{frame}

\begin{frame}{標準化係数 (Beta Coefficients)}
\protect\hypertarget{ux6a19ux6e96ux5316ux4fc2ux6570-beta-coefficients}{}
\begin{itemize}
\tightlist
\item
  変数の単位に依存しない係数比較を行いたい場合、変数を\textbf{標準化}（平均を引いて標準偏差で割る）する。
\item
  すべての変数を標準化して回帰を行うと、得られる係数は\textbf{標準化係数}（Beta
  Coefficients）と呼ばれる。
\item
  「\(x\) が1標準偏差変化したとき、\(y\)
  が何標準偏差変化するか」を表す。
\end{itemize}
\end{frame}

\hypertarget{ux95a2ux6570ux5f62ux5f0fux306eux3055ux3089ux306aux308bux691cux8a0e}{%
\section{6-2
関数形式のさらなる検討}\label{ux95a2ux6570ux5f62ux5f0fux306eux3055ux3089ux306aux308bux691cux8a0e}}

\begin{frame}{対数形式 (Log Forms)}
\protect\hypertarget{ux5bfeux6570ux5f62ux5f0f-log-forms}{}
\begin{itemize}
\tightlist
\item
  対数形式をとることで、非線形な関係を線形モデルで扱える。
\item
  対数の差は近似的に変化率（パーセント）を表す。
\end{itemize}

\begin{longtable}[]{@{}llll@{}}
\toprule
モデル & 従属変数 & 独立変数 & 解釈 (\(\beta_1\)) \\
\midrule
\endhead
Level-Level & \(y\) & \(x\) & \(\Delta y = \beta_1 \Delta x\) \\
Level-Log & \(y\) & \(\log(x)\) &
\(\Delta y \approx (\beta_1/100) \% \Delta x\) \\
Log-Level & \(\log(y)\) & \(x\) &
\(\% \Delta y \approx (100 \beta_1) \Delta x\) \\
Log-Log & \(\log(y)\) & \(\log(x)\) &
\(\% \Delta y \approx \beta_1 \% \Delta x\) (弾力性) \\
\bottomrule
\end{longtable}
\end{frame}

\begin{frame}{対数形式の利点}
\protect\hypertarget{ux5bfeux6570ux5f62ux5f0fux306eux5229ux70b9}{}
\begin{enumerate}
\tightlist
\item
  係数が変化率（弾力性）として解釈しやすくなる。
\item
  変数の分布が右に裾を引いている場合（所得など）、対数をとることで正規分布に近づき、外れ値の影響を緩和できる。
\item
  不均一分散を緩和することが多い。
\end{enumerate}
\end{frame}

\begin{frame}{2次形式 (Quadratic Forms)}
\protect\hypertarget{ux6b21ux5f62ux5f0f-quadratic-forms}{}
\begin{itemize}
\tightlist
\item
  限界効果が一定でない場合に使用。
\item
  モデル: \(y = \beta_0 + \beta_1 x + \beta_2 x^2 + u\)
\item
  限界効果: \(\frac{\partial y}{\partial x} = \beta_1 + 2\beta_2 x\)
\item
  \(\beta_2 > 0\): 下に凸（U字型）
\item
  \(\beta_2 < 0\): 上に凸（逆U字型）
\item
  頂点（極値）: \(x^* = -\beta_1 / (2\beta_2)\)
\end{itemize}
\end{frame}

\begin{frame}[fragile]{2次形式の例 (\texttt{hprice1})}
\protect\hypertarget{ux6b21ux5f62ux5f0fux306eux4f8b-hprice1}{}
\footnotesize

\begin{Shaded}
\begin{Highlighting}[]
\FunctionTok{data}\NormalTok{(hprice1)}
\CommentTok{\# 住宅面積(sqrft)の2次形式}
\NormalTok{res }\OtherTok{\textless{}{-}} \FunctionTok{lm}\NormalTok{(price }\SpecialCharTok{\textasciitilde{}}\NormalTok{ lotsize }\SpecialCharTok{+}\NormalTok{ sqrft }\SpecialCharTok{+} \FunctionTok{I}\NormalTok{(sqrft}\SpecialCharTok{\^{}}\DecValTok{2}\NormalTok{) }\SpecialCharTok{+}\NormalTok{ bdrms, }\AttributeTok{data=}\NormalTok{hprice1)}
\FunctionTok{summary}\NormalTok{(res)}\SpecialCharTok{$}\NormalTok{coefficients[}\DecValTok{3}\SpecialCharTok{:}\DecValTok{4}\NormalTok{, ]}
\end{Highlighting}
\end{Shaded}

\begin{verbatim}
##                 Estimate   Std. Error    t value   Pr(>|t|)
## sqrft      -8.622611e-03 6.953558e-02 -0.1240029 0.90161283
## I(sqrft^2)  2.741764e-05 1.425207e-05  1.9237662 0.05781174
\end{verbatim}

\normalsize
\end{frame}

\begin{frame}{交差項 (Interaction Terms)}
\protect\hypertarget{ux4ea4ux5deeux9805-interaction-terms}{}
\begin{itemize}
\tightlist
\item
  ある独立変数の効果が、別の独立変数の値に依存する場合。
\item
  モデル:
  \(y = \beta_0 + \beta_1 x_1 + \beta_2 x_2 + \beta_3 (x_1 x_2) + u\)
\item
  \(x_1\) の限界効果:
  \(\frac{\partial y}{\partial x_1} = \beta_1 + \beta_3 x_2\)
\item
  \(x_2\) の値によって \(x_1\) の効果が変化することを表現できる。
\end{itemize}
\end{frame}

\hypertarget{ux6c7aux5b9aux4fc2ux6570-r2-ux3068ux81eaux7531ux5ea6ux4feeux6b63ux6c7aux5b9aux4fc2ux6570-barr2}{%
\section{\texorpdfstring{6-3 決定係数 \(R^2\) と自由度修正決定係数
\(\bar{R}^2\)}{6-3 決定係数 R\^{}2 と自由度修正決定係数 \textbackslash bar\{R\}\^{}2}}\label{ux6c7aux5b9aux4fc2ux6570-r2-ux3068ux81eaux7531ux5ea6ux4feeux6b63ux6c7aux5b9aux4fc2ux6570-barr2}}

\begin{frame}{通常の \(R^2\) の問題点}
\protect\hypertarget{ux901aux5e38ux306e-r2-ux306eux554fux984cux70b9}{}
\begin{itemize}
\tightlist
\item
  モデルに独立変数を追加すると、たとえその変数に説明力がなくても、\(R^2\)
  は決して減少せず、通常は増加する。
\item
  モデルの複雑さを考慮した指標が必要。
\end{itemize}
\end{frame}

\begin{frame}{自由度修正決定係数 (Adjusted R-squared)}
\protect\hypertarget{ux81eaux7531ux5ea6ux4feeux6b63ux6c7aux5b9aux4fc2ux6570-adjusted-r-squared}{}
\begin{itemize}
\tightlist
\item
  推定に要した自由度を考慮した決定係数：
  \[ \bar{R}^2 = 1 - \frac{SSR / (n - k - 1)}{SST / (n - 1)} = 1 - \frac{\hat{\sigma}^2}{SST / (n-1)} \]
\item
  変数を追加しても、\(\hat{\sigma}^2\)（誤差分散の推定値）が減少しなければ
  \(\bar{R}^2\) は増加しない。
\item
  \(\bar{R}^2\) は負になることもある。
\end{itemize}
\end{frame}

\begin{frame}{モデル選択における注意}
\protect\hypertarget{ux30e2ux30c7ux30ebux9078ux629eux306bux304aux3051ux308bux6ce8ux610f}{}
\begin{itemize}
\tightlist
\item
  \(R^2\) や \(\bar{R}^2\)
  を最大化することだけが、モデル選択の目的ではない。
\item
  重要なのは、理論的背景や、内生性（\(E[u|x]=0\)）の排除である。
\item
  非線形な変換（例：\(y\) vs \(\log(y)\)）をしたモデル同士の \(R^2\)
  を直接比較してはいけない。
\end{itemize}
\end{frame}

\hypertarget{ux4e88ux6e2cux3068ux6b8bux5deeux5206ux6790}{%
\section{6-4
予測と残差分析}\label{ux4e88ux6e2cux3068ux6b8bux5deeux5206ux6790}}

\begin{frame}[fragile]{予測 (Prediction)}
\protect\hypertarget{ux4e88ux6e2c-prediction}{}
\begin{itemize}
\tightlist
\item
  推定されたモデルを用いて、与えられた \(x\) の値に対する \(y\)
  の期待値を推定する。
\item
  \texttt{predict()} 関数を使用して計算可能。
\item
  信頼区間 (Confidence Interval) と予測区間 (Prediction Interval)
  の違いに注意。
\end{itemize}
\end{frame}

\begin{frame}{残差分析 (Residual Analysis)}
\protect\hypertarget{ux6b8bux5deeux5206ux6790-residual-analysis}{}
\begin{itemize}
\tightlist
\item
  推定された残差 \(\hat{u} = y - \hat{y}\)
  を観察することで、モデルの妥当性をチェックする。
\item
  特定のサンプル（アウトライヤー）が結果に大きな影響を与えていないか確認する。
\end{itemize}
\end{frame}

\begin{frame}{まとめ}
\protect\hypertarget{ux307eux3068ux3081}{}
\begin{itemize}
\tightlist
\item
  スケーリングは係数の値を変化させるが、t統計量や \(R^2\)
  には影響しない。
\item
  対数、2次形式、交差項により、現実的な非線形関係をモデル化できる。
\item
  \(\bar{R}^2\) はモデルの複雑さを考慮するが、過信は禁物。
\item
  予測と残差分析は、モデルの実用性と妥当性の検証に不可欠。
\end{itemize}
\end{frame}

\end{document}
